\documentclass[11pt]{article}
\usepackage[margin=1in]{geometry}
\usepackage{amsmath,amssymb,amsthm}
\usepackage{graphicx}
\usepackage[hidelinks]{hyperref}

\newtheorem{theorem}{Theorem}
\newtheorem{lemma}{Lemma}
\newtheorem{corollary}{Corollary}
\newtheorem{definition}{Definition}

\newcommand{\N}{\mathbb N}
\newcommand{\eps}{\varepsilon}
\newcommand{\linf}{\ell^\infty}
\newcommand{\Id}{I}
\newcommand{\cD}{\mathcal D}

\title{A Recurrent-Exponent Self-Similarity Criterion for\\
the \texorpdfstring{$\ell^\infty(L^{p_k})$}{ell-infinity(Lp_k)} Factorization Problem}
\author{Run \texttt{fa\_banach\_001}, lane 8}
\date{2026-06-23}

\begin{document}
\maketitle

\section*{Status}

This is a partial result for the final questions of Lechner, Motakis, Mueller, and
Schlumprecht~\cite{LMMMS}.  It proves a broad self-similar mixed-exponent case of the
large-diagonal factorization problem for the Haar array in
\[
  Z(p)=\linf(L^{p_k}:k\in\N).
\]
It does not settle the strict endpoint cases \(p_k\to 1\) or \(p_k\to\infty\) from Question~2 of
the source paper.

\section*{Source Question}

The source paper asks the following on page 26.

\begin{center}
\includegraphics[width=0.95\linewidth]{figures/open_problem_crop.png}
\end{center}

In text form:
\begin{quote}
Assume that \((p_k)\), \(1\leq p_k<\infty\), either converges to \(1\) or diverges to
\(\infty\). Does the array \((h_{k,I})\) have the factorization property in
\(Z=\ell^\infty(L^{p_k}:k\in\mathbb N)\)?
\end{quote}

The present packet proves a nearby structural criterion that recovers the dense-exponent example in
the source paper and extends it to arbitrary recurrent exponent lists.

\section*{Definitions}

For a sequence \(p=(p_k)\) in \([1,\infty)\), write
\[
  Z(p)=\linf(L^{p_k}:k\in\N).
\]
The Haar array in \(Z(p)\) is denoted by \((h_{k,I})_{k\in\N,I\in\cD}\), where \(h_{k,I}\) is the
usual \(L^\infty\)-normalized Haar function \(h_I\) placed in the \(k\)-th coordinate.

\begin{definition}
An exponent sequence \(p=(p_k)\subset [1,\infty)\) is \emph{recurrent} if, for every \(k\in\N\) and
every \(\eps>0\), the set
\[
  \{m\in\N: |p_m-p_k|<\eps\}
\]
is infinite.
\end{definition}

\section*{Proof Intuition}

The source paper proves a powerful subsequence theorem: under the Haar hypotheses, every operator
on \(Z(p)\) with large diagonal factors the identity on a sub-sum \(Z(p|_\Gamma)\), and \(\Gamma\)
can be forced to meet any prescribed countable family of infinite sets.  To turn this into a
factorization of \(\Id_{Z(p)}\), one needs \(Z(p|_\Gamma)\) to be another copy of \(Z(p)\).

Recurrence is exactly a clean way to guarantee this.  It lets us prescribe infinitely many future
coordinates whose exponents approximate each original exponent \(p_k\) to the finite-dimensional
accuracy needed in Proposition 7.2 of the source paper.  The Johnson-ultrafilter construction used
there then embeds each \(L^{p_k}\) complementably into the selected sub-sum.  Running this
coordinatewise and applying Pelczynski decomposition gives \(Z(p|_\Gamma)\cong Z(p)\).  Composing
this isomorphism with the source factorization on \(Z(p|_\Gamma)\) gives the desired factorization
of \(\Id_{Z(p)}\).

\section*{Main Result}

\begin{theorem}
\label{thm:main}
Let \(p=(p_k)\subset[1,\infty)\) be recurrent.  Then the Haar array
\((h_{k,I})_{k,I}\) has the factorization property in \(Z(p)=\ell^\infty(L^{p_k}:k\in\N)\).  In
other words, if \(T:Z(p)\to Z(p)\) is bounded and
\[
  \inf_{k,I} |h_{k,I}^*(T h_{k,I})|>0,
\]
then \(\Id_{Z(p)}\) factors through \(T\).
\end{theorem}

\begin{lemma}
\label{lem:cluster-isomorphism}
Let \(p=(p_k)\) and \(q=(q_k)\) be sequences in \([1,\infty)\).  Assume:
\begin{enumerate}
\item for every \(k\), \(p_k\) is an accumulation value of the sequence \(q\);
\item for every \(k\), \(q_k\) is an accumulation value of the sequence \(p\);
\item \(Z(p)\) is isomorphic to \(\ell^\infty(Z(p))\).
\end{enumerate}
Then \(Z(q)\) is isomorphic to \(Z(p)\).
\end{lemma}

\begin{proof}
We first recall the argument of Proposition 7.2 in~\cite{LMMMS}.  If a fixed exponent \(r\) is an
accumulation value of \(q\), choose pairwise distinct indices \(m_s\) such that \(q_{m_s}\to r\) and
so fast that the identity map between the first \(s\) Haar levels of \(L^{q_{m_s}}\) and
\(L^r\) is a \(2\)-isomorphism.  The map
\[
  f\in L^r \longmapsto (R_s f)_{s=1}^\infty\in\ell^\infty(L^{q_{m_s}})
\]
where \(R_s\) is the \(s\)-th Haar partial-sum projection, is then an isomorphic embedding.  The
ultrafilter coefficient-limit map used in Proposition 7.2 gives a bounded left inverse, so the
image is complemented.  The constants are uniform after fixing the numerical tolerance \(2\).

Applying this construction for each coordinate \(p_k\), and using disjoint subsequences of the
\(q\)-coordinates, gives a complemented embedding of \(Z(p)\) into \(Z(q)\).  The same argument with
the roles of \(p\) and \(q\) interchanged gives a complemented embedding of \(Z(q)\) into \(Z(p)\).
Since \(Z(p)\cong\ell^\infty(Z(p))\), Pelczynski decomposition yields \(Z(q)\cong Z(p)\).
\end{proof}

\begin{lemma}
\label{lem:recurrent-accordion}
If \(p\) is recurrent, then \(Z(p)\cong\ell^\infty(Z(p))\).
\end{lemma}

\begin{proof}
Let \(p^\infty\) denote the sequence obtained by repeating every \(p_k\) countably many times.  Then
\(Z(p^\infty)\) is isometric to \(\ell^\infty(Z(p))\).  The sequence \(p^\infty\) trivially has each
\(p_k\) as an accumulation value, and recurrence of \(p\) lets us allocate disjoint subsequences of
\(p\) approximating each countably repeated exponent in \(p^\infty\).  Thus the proof of
Lemma~\ref{lem:cluster-isomorphism}, applied to \(p\) and \(p^\infty\), gives complemented
embeddings both ways.  Since \(Z(p^\infty)\) is visibly isomorphic to its own \(\ell^\infty\)-sum,
Pelczynski decomposition gives \(Z(p)\cong Z(p^\infty)\cong \ell^\infty(Z(p))\).
\end{proof}

\begin{proof}[Proof of Theorem~\ref{thm:main}]
For each pair \((k,s)\), choose \(\eta_{k,s}>0\) so small that whenever
\(|r-p_k|<\eta_{k,s}\), the coefficient identity between the first \(s\) Haar levels of \(L^{r}\)
and \(L^{p_k}\) is a \(2\)-isomorphism.  Put
\[
  \Omega_{k,s}=\{m\in\N: |p_m-p_k|<\eta_{k,s}\}.
\]
By recurrence, every \(\Omega_{k,s}\) is infinite.

Theorem 2.5 of~\cite{LMMMS} applies to the Haar array in \(Z(p)\): the Haar bases have uniformly
bounded basis constants, the required diagonal factorization property, and simultaneous strategic
reproducibility, as recorded in the source paper for the \(L^p\) Haar systems.  Therefore, for any
operator \(T:Z(p)\to Z(p)\) with large diagonal, there is an infinite
\(\Gamma\subset\N\) such that \(\Gamma\cap\Omega_{k,s}\) is infinite for every \(k,s\), and
\[
  \Id_{Z(p|_\Gamma)}
\]
factors through \(T\).

The intersection property implies that every exponent \(p_k\) is an accumulation value of the
subsequence \(p|_\Gamma\).  Conversely, every exponent in \(p|_\Gamma\) is one of the original
\(p_k\)'s, hence is an accumulation value of \(p\) by recurrence.  By
Lemmas~\ref{lem:cluster-isomorphism} and~\ref{lem:recurrent-accordion},
\[
  Z(p|_\Gamma)\cong Z(p).
\]
Composing this isomorphism with the source factorization of \(\Id_{Z(p|_\Gamma)}\) through \(T\)
yields a factorization of \(\Id_{Z(p)}\) through \(T\).  This proves the theorem.
\end{proof}

\section*{Limitations}

The theorem does not resolve Question 2 as stated.  A strict sequence \(p_k\to1\), or a strict
sequence \(p_k\to\infty\), usually has many isolated exponents: for a fixed early \(p_k\), there may
be no infinite tail of exponents approximating \(p_k\).  The argument above then cannot prove that
the selected sub-sum \(Z(p|_\Gamma)\) is isomorphic to the original \(Z(p)\).  Equivalently, the
remaining endpoint obstruction is an absorption problem for isolated \(L^p\)-coordinates inside
tails whose exponents converge to \(1\) or to \(\infty\).

\section*{Novelty and Literature Check}

The source paper itself proves the dense-exponent case by showing that if \((p_k)\) is dense in
\([1,\infty)\), then the selected sub-sum remains isomorphic to \(Z(p)\).  The present packet
abstracts that mechanism to a recurrence/self-similarity criterion.  Local checks of later
factorization-related sources in the run cache, including arXiv:2011.09915, 2102.10088, 2112.12534,
and 2605.21711, found related factorization, primariness, and Haar-system results but no direct
answer to the endpoint question and no exact statement of this recurrent-exponent criterion.
Web searches for phrases around arXiv:1910.11188, \(\ell^\infty(L^{p_k})\), Haar arrays, and the
endpoint Question 2 also found no direct answer beyond the source paper itself.

\section*{Human Review Recommendation}

Review as a likely valid partial result.  The key points to check are:
\begin{enumerate}
\item the adaptation of Proposition 7.2 of~\cite{LMMMS} from dense exponent sequences to arbitrary
cluster sequences;
\item the coordinatewise complemented-embedding construction using disjoint approximating
subsequences;
\item the Pelczynski decomposition step converting mutual complemented embeddings into an
isomorphism.
\end{enumerate}

\begin{thebibliography}{9}
\bibitem{LMMMS}
R. Lechner, P. Motakis, P. F. X. Mueller, and Th. Schlumprecht,
\emph{The factorization property of \(\ell^\infty(X_k)\)}, arXiv:1910.11188, 2019.
\end{thebibliography}

\end{document}
